\documentclass[10pt,a4paper]{amsart}
\usepackage{amsmath,amssymb,amsthm,graphicx}
% For commutative diagrams we use Xy-pic
\usepackage[all]{xy}

% We use 2cell for 2-commutative diagrams.
\xyoption{2cell}
% Theorem environments.
%
\theoremstyle{plain}
\newtheorem{theorem}[subsection]{Theorem}
\newtheorem{proposition}[subsection]{Proposition}
\newtheorem{lemma}[subsection]{Lemma}

\theoremstyle{definition}
\newtheorem{definition}[subsection]{Definition}
\newtheorem{example}[subsection]{Example}
\newtheorem{exercise}[subsection]{Exercise}
\newtheorem{situation}[subsection]{Situation}

\theoremstyle{remark}
\newtheorem{remark}[subsection]{Remark}
\newtheorem{remarks}[subsection]{Remarks}

\numberwithin{equation}{subsection}

% Macros
%
\def\lim{\mathop{\mathrm{lim}}\nolimits}
\def\colim{\mathop{\mathrm{colim}}\nolimits}
\def\Spec{\mathop{\mathrm{Spec}}}
\def\Hom{\mathop{\mathrm{Hom}}\nolimits}
\def\Ext{\mathop{\mathrm{Ext}}\nolimits}
\def\SheafHom{\mathop{\mathcal{H}\!\mathit{om}}\nolimits}
\def\SheafExt{\mathop{\mathcal{E}\!\mathit{xt}}\nolimits}
\def\Sch{\mathit{Sch}}
\def\Mor{\mathop{\mathrm{Mor}}\nolimits}
\def\Ob{\mathop{\mathrm{Ob}}\nolimits}
\def\Sh{\mathop{\mathit{Sh}}\nolimits}
\def\NL{\mathop{N\!L}\nolimits}
\def\CH{\mathop{\mathrm{CH}}\nolimits}
\def\proetale{{pro\text{-}\acute{e}tale}}
\def\etale{{\acute{e}tale}}
\def\QCoh{\mathit{QCoh}}
\def\Ker{\mathop{\mathrm{Ker}}}
\def\Im{\mathop{\mathrm{Im}}}
\def\Coker{\mathop{\mathrm{Coker}}}
\def\Coim{\mathop{\mathrm{Coim}}}

% Boxtimes
%
\DeclareMathSymbol{\boxtimes}{\mathbin}{AMSa}{"02}

%
% Macros for moduli stacks/spaces
%
\def\QCohstack{\mathcal{QC}\!\mathit{oh}}
\def\Cohstack{\mathcal{C}\!\mathit{oh}}
\def\Spacesstack{\mathcal{S}\!\mathit{paces}}
\def\Quotfunctor{\mathrm{Quot}}
\def\Hilbfunctor{\mathrm{Hilb}}
\def\Curvesstack{\mathcal{C}\!\mathit{urves}}
\def\Polarizedstack{\mathcal{P}\!\mathit{olarized}}
\def\Complexesstack{\mathcal{C}\!\mathit{omplexes}}
% \Pic is the operator that assigns to X its picard group, usage \Pic(X)
% \Picardstack_{X/B} denotes the Picard stack of X over B
% \Picardfunctor_{X/B} denotes the Picard functor of X over B
\def\Pic{\mathop{\mathrm{Pic}}\nolimits}
\def\Picardstack{\mathcal{P}\!\mathit{ic}}
\def\Picardfunctor{\mathrm{Pic}}
\def\Deformationcategory{\mathcal{D}\!\mathit{ef}}
\usepackage{cleveref}
\crefname{originalreference}{original reference}{original references}
\Crefname{originalreference}{Original reference}{Original references}
\makeatletter
\newlabel{lemma-compatible-shriek-push-finite}{{0GJ1}{source}{}{sourcepointer-1}{}}
\newlabel{lemma-compatible-shriek-push-finite@cref}{{[originalreference][0][]0GJ1}{[1][0][]source}{}{}{}}
\makeatother
\crefname{originalreference}{original reference}{original references}
\Crefname{originalreference}{Original reference}{Original references}
\makeatletter
\newlabel{lemma-finite-pushforward-coefficients}{{0GJD}{source}{}{sourcepointer-2}{}}
\newlabel{lemma-finite-pushforward-coefficients@cref}{{[originalreference][0][]0GJD}{[1][0][]source}{}{}{}}
\makeatother
\crefname{originalreference}{original reference}{original references}
\Crefname{originalreference}{Original reference}{Original references}
\makeatletter
\newlabel{lemma-pullback-filtered-modules}{{0GJ0}{source}{}{sourcepointer-3}{}}
\newlabel{lemma-pullback-filtered-modules@cref}{{[originalreference][0][]0GJ0}{[1][0][]source}{}{}{}}
\makeatother
\crefname{originalreference}{original reference}{original references}
\Crefname{originalreference}{Original reference}{Original references}
\makeatletter
\newlabel{lemma-ses-statements-coefficients}{{0GJG}{source}{}{sourcepointer-4}{}}
\newlabel{lemma-ses-statements-coefficients@cref}{{[originalreference][0][]0GJG}{[1][0][]source}{}{}{}}
\makeatother
\crefname{originalreference}{original reference}{original references}
\Crefname{originalreference}{Original reference}{Original references}
\makeatletter
\newlabel{lemma-somewhat-easier-coefficients}{{0GJF}{source}{}{sourcepointer-5}{}}
\newlabel{lemma-somewhat-easier-coefficients@cref}{{[originalreference][0][]0GJF}{[1][0][]source}{}{}{}}
\makeatother
\crefname{originallemma}{original lemma}{original lemmas}
\Crefname{originallemma}{Original lemma}{Original lemmas}
\makeatletter
\newlabel{more-morphisms-lemma-quasi-finite-separated-pass-through-finite}{{05K0}{source}{}{sourcepointer-6}{}}
\newlabel{more-morphisms-lemma-quasi-finite-separated-pass-through-finite@cref}{{[originallemma][0][]05K0}{[1][0][]source}{}{}{}}
\makeatother
\crefname{originalsection}{original section}{original sections}
\Crefname{originalsection}{Original section}{Original sections}
\makeatletter
\newlabel{section-trace-method}{{03SH}{source}{}{sourcepointer-7}{}}
\newlabel{section-trace-method@cref}{{[originalsection][0][]03SH}{[1][0][]source}{}{}{}}
\makeatother
\title{Finiteness of the cohomology of the extension by zero of a locally constant sheaf annihilated by a prime}
\author{The Stacks Project Authors}
\date{stacks-project:etale-cohomology.tex:local-snapshot-20261011; GNU FDL 1.2 or later}
\begin{document}
\maketitle
\noindent\textit{Editorial note.} All delivered mathematical text is the author's own, copied from the pinned Stacks Project chapter file etale-cohomology.tex; no statement or proof step was rewritten, reordered or removed.\par\medskip
\noindent\textit{Editorial note.} The counted proof is the author's complete argument: the prime-to-l finite etale cover and the trace-method splitting step, the Zariski-factorization step, the shriek/pushforward identity used to move the extension by zero across the finite morphism, and the filtration-plus-short-exact-sequence reduction to the constant case.\par\medskip
\noindent\textit{Editorial note.} The label lemma-vanishing-easier-coefficients names the counted statement, and the proof environment immediately following it in the source is delivered verbatim as its proof.\par\medskip
\noindent\textit{Editorial note.} The counted statement is carried by the author's own lemma environment in the chapter source.\par\medskip
\noindent\textit{Editorial note.} Pre-open gate: python operations/agent-authoring/tools/precard\_gate.py --source etale-cohomology.tex --label lemma-vanishing-easier-coefficients --start 16010 --end 16059 --local-id 3553 returned RESULT=OK (L3553 free, occupancy FREE, public interval clear). Candidate taken from lane-s4/halfB-priority.json.\par\medskip
\noindent\textit{Editorial note.} The statement and proof are copied from the pinned chapter file; the cited inputs are kept as external pointers with their canonical Stacks tags.\par\medskip
\noindent\textit{Allowed standard prerequisites (editorial).}\begin{itemize}
\item The trace method and the finite etale cover of degree prime to l. Same chapter, the section on the trace method and the lemma producing a finite etale V -> U of degree prime to l together with a filtration of the pullback. Supplies the cover f : V -> U of degree prime to l and the maps G -> f\_*f\textasciicircum{}\{-1\}G -> G whose composition is an isomorphism, and the resulting filtration of the pulled-back coefficient sheaf.
\item Zariski's main theorem, quasi-finite separated version. More on Morphisms, the lemma that a quasi-finite separated morphism factors as an open immersion followed by a finite morphism. Supplies the diagram with Y -> X finite and j' : V -> Y an open immersion with dense image.
\item Shriek pushforward versus pushforward for a finite morphism. Same chapter, the lemma identifying j\_! composed with the pushforward along a finite morphism with the pushforward of the shriek extension. Supplies the identity j\_!f\_*f\textasciicircum{}\{-1\}G = fbar\_*j'\_!f\textasciicircum{}\{-1\}G used to replace the extension by zero along j by an extension by zero along j'.
\item Pushforward of finite type coefficient sheaves along a finite morphism. Same chapter, the lemma reducing finiteness of the cohomology of a pushforward along a finite morphism to the source. Supplies the reduction from j'\_!f\textasciicircum{}\{-1\}G on Y to the sheaf j'\_!f\textasciicircum{}\{-1\}G on Y with the same finiteness conclusions.
\item Finiteness for a finite constant coefficient sheaf, and the short exact sequence reduction. Same chapter, the lemma treating the case of a constant finite coefficient sheaf and the lemma reducing along a short exact sequence of coefficient sheaves. Supplies the final reduction from a filtered coefficient sheaf with finite quotients to the constant case M.
\end{itemize}\medskip
\ifdefined\hypertarget\hypertarget{sourcepointer-1}{}\fi
\noindent\textit{Source reference.} Original reference 0GJ1: sources/stacks-etale-cohomology/pinned-original-source.tex, source lines 11612--11612. Prerequisite pointer invoked at proof line(s) 16047; it names a standard fact of the same chapter that the counted argument uses, and every step of the counted proof built on it is retained verbatim in the delivered body.\par\medskip
\ifdefined\hypertarget\hypertarget{sourcepointer-2}{}\fi
\noindent\textit{Source reference.} Original reference 0GJD: sources/stacks-etale-cohomology/pinned-original-source.tex, source lines 15909--15909. Prerequisite pointer invoked at proof line(s) 16051; it names a standard fact of the same chapter that the counted argument uses, and every step of the counted proof built on it is retained verbatim in the delivered body.\par\medskip
\ifdefined\hypertarget\hypertarget{sourcepointer-3}{}\fi
\noindent\textit{Source reference.} Original reference 0GJ0: sources/stacks-etale-cohomology/pinned-original-source.tex, source lines 10479--10479. Prerequisite pointer invoked at proof line(s) 16028; it names a standard fact of the same chapter that the counted argument uses, and every step of the counted proof built on it is retained verbatim in the delivered body.\par\medskip
\ifdefined\hypertarget\hypertarget{sourcepointer-4}{}\fi
\noindent\textit{Source reference.} Original reference 0GJG: sources/stacks-etale-cohomology/pinned-original-source.tex, source lines 15997--15997. Prerequisite pointer invoked at proof line(s) 16055; it names a standard fact of the same chapter that the counted argument uses, and every step of the counted proof built on it is retained verbatim in the delivered body.\par\medskip
\ifdefined\hypertarget\hypertarget{sourcepointer-5}{}\fi
\noindent\textit{Source reference.} Original reference 0GJF: sources/stacks-etale-cohomology/pinned-original-source.tex, source lines 15973--15973. Prerequisite pointer invoked at proof line(s) 16058; it names a standard fact of the same chapter that the counted argument uses, and every step of the counted proof built on it is retained verbatim in the delivered body.\par\medskip
\ifdefined\hypertarget\hypertarget{sourcepointer-6}{}\fi
\noindent\textit{Source reference.} Original lemma 05K0: sources/stacks-chow/tags.txt, source lines 7184--7184. Prerequisite pointer invoked at proof line(s) 16037; it names a standard fact of a sibling chapter that the counted argument uses, and every step of the counted proof built on it is retained verbatim in the delivered body.\par\medskip
\ifdefined\hypertarget\hypertarget{sourcepointer-7}{}\fi
\noindent\textit{Source reference.} Original section 03SH: sources/stacks-etale-cohomology/pinned-original-source.tex, source lines 10330--10330. Prerequisite pointer invoked at proof line(s) 16029; it names a standard fact of the same chapter that the counted argument uses, and every step of the counted proof built on it is retained verbatim in the delivered body.\par\medskip
\begin{lemma}
\label{lemma-vanishing-easier-coefficients}
Let $\Lambda$ be a Noetherian ring, let $k$ be an algebraically closed field,
let $X$ be a separated, finite type scheme of dimension $\leq 1$ over $k$,
let $j : U \to X$ be an open immersion with $U$ connected, let
$\ell$ be a prime number, let $n > 0$, and let $\mathcal{G}$ be a finite type,
locally constant sheaf of $\Lambda$-modules on $U_\etale$ annihilated by
$\ell^n$. Then
\begin{enumerate}
\item $H^q_\etale(X, j_!\mathcal{G})$ is a finite $\Lambda$-module
if $\ell$ is prime to $\text{char}(k)$,
\item $H^q_\etale(X, j_!\mathcal{G})$ is a finite $\Lambda$-module
if $X$ is proper.
\end{enumerate}
\end{lemma}

\begin{proof}
Let $f : V \to U$ be a finite \'etale morphism of degree prime to $\ell$
as in Lemma \ref{lemma-pullback-filtered-modules}. The discussion in
Section \ref{section-trace-method} gives maps
$$
\mathcal{G} \to f_*f^{-1}\mathcal{G} \to \mathcal{G}
$$
whose composition is an isomorphism. Hence it suffices to prove the
finiteness of $H^q_\etale(X, j_!f_*f^{-1}\mathcal{G})$.
By Zariski's Main theorem
(More on Morphisms, Lemma
\ref{more-morphisms-lemma-quasi-finite-separated-pass-through-finite})
we can choose a diagram
$$
\xymatrix{
V \ar[r]_{j'} \ar[d]_f & Y \ar[d]^{\overline{f}} \\
U \ar[r]^j & X
}
$$
with $\overline{f} : Y \to X$ finite and $j'$ an open immersion with
dense image. Since $f$ is finite and $V$ dense in $Y$ we have
$V = U \times_X Y$. By Lemma \ref{lemma-compatible-shriek-push-finite} we have
$$
j_!f_*f^{-1}\mathcal{G} = \overline{f}_*j'_!f^{-1}\mathcal{G}
$$
By Lemma \ref{lemma-finite-pushforward-coefficients} it suffices to
consider $j'_!f^{-1}\mathcal{G}$. The existence of the filtration given by
Lemma \ref{lemma-pullback-filtered-modules},
the fact that $j'_!$ is exact, and
Lemma \ref{lemma-ses-statements-coefficients}
reduces us to the case
$\mathcal{F} = j'_!\underline{M}$ for a finite $\Lambda$-module $M$
which is Lemma \ref{lemma-somewhat-easier-coefficients}.
\end{proof}

\end{document}
